\documentclass[11pt]{article}
\usepackage[T1]{fontenc}
\usepackage{textcomp}
\usepackage[margin=0.75in]{geometry}
\usepackage{amsmath,amssymb,amsthm}
\usepackage{array,booktabs}
\usepackage{hyperref}
\setlength{\parindent}{0pt}
\newtheorem{theorem}{Theorem}
\theoremstyle{definition}
\newtheorem{definition}{Definition}
\title{Flow Proof Artifact: thales-central-inscribed-witness}
\author{Generated by \texttt{flow doc proof}}
\date{Flow Proof Book}
\begin{document}
\maketitle
Central and inscribed angles on a diameter satisfy the half relation.\\[0.5em]
\textbf{Source.} euclid — Elements, Book III, Propositions 20 and 31\\[0.5em]
\bigskip
\noindent{\large \textbf{Derived fact 1.} \textit{angle ACB equals half of angle AOB on diameter AB} \hfill the Euclidean plane \cdot circle \cdot central inscribed diameter witness}\par
\smallskip\noindent\textit{Coordinate: the Euclidean plane \cdot circle \cdot central inscribed diameter witness}\par
\smallskip\noindent\textit{Source: euclid}\par
\smallskip\noindent\textit{Built on: inscribed central ratio is half, for circles in on the Euclidean plane, inscribed on semicircle is right, for circles in on the Euclidean plane}\par
\medskip
\noindent\textbf{Goal.} angle ACB equals half of angle AOB on diameter AB\par
\noindent$\displaystyle \angle ACB = \tfrac{1}{2}\angle AOB$\par
\medskip
\noindent
\begin{tabular}{@{} >{\raggedright\arraybackslash}p{0.06\textwidth} >{\raggedright\arraybackslash}p{0.47\textwidth} >{\raggedleft\arraybackslash}p{0.41\textwidth} @{}}
\textbf{\#} & \textbf{Proof} & \textbf{Mathematics} \\
\midrule
\textbf{1.} & We prove that central inscribed diameter witness for circles in the Euclidean plane. & \\[0.45em]
\textbf{2.} & We invoke the derived fact governing circles in the Euclidean plane: inscribed central ratio is half, for circles in on the Euclidean plane. & \\[0.45em]
\textbf{3.} & We invoke the derived fact governing circles in the Euclidean plane: inscribed on semicircle is right, for circles in on the Euclidean plane. & \\[0.45em]
\textbf{4.} & From step 2 and step 3, this implies angle ACB equals half of angle AOB. Hence proven. & $\displaystyle \angle ACB = \tfrac{1}{2}\angle AOB$ \\[0.45em]
\bottomrule
\end{tabular}
\label{thm:Geometry-circle-central_inscribed_diameter_witness}
\par\medskip\hrule
\end{document}
